\IfFileExists{stacks-project.cls}{%
\documentclass{stacks-project}
}{%
\documentclass{amsart}
}
\usepackage{amssymb}

% For dealing with references we use the comment environment
\usepackage{verbatim}
\usepackage{fancyvrb}
\usepackage{pmboxdraw}
\newenvironment{reference}{\comment}{\endcomment}
%\newenvironment{reference}{}{}
\newenvironment{slogan}{\comment}{\endcomment}
\newenvironment{history}{\comment}{\endcomment}

% For source figures retained as assets
\usepackage{graphicx}

% For commutative diagrams we use Xy-pic
\usepackage[all]{xy}

% We use 2cell for 2-commutative diagrams.
\xyoption{2cell}
\UseAllTwocells

% We use multicol for the list of chapters between chapters
\usepackage{multicol}

% This is generally recommended for better output
\usepackage{lmodern}
\usepackage[T1]{fontenc}
\usepackage{textalpha}

% For cross-file-references
% We prefer the snapshot tmp/xref/#2.aux,
% which latexmkrc replaces atomically, so
% that documents compiled simultaneously
% never read a half-written .aux file.
% The last argument keeps links to #2.pdf.
\usepackage{xr-hyper}
\newcommand{\maybeexternaldocument}[2]{%
  \IfFileExists{#2.aux}{%
    \externaldocument[#1]{#2}%
  }{%
    \PackageWarning{stack}{Missing #2.aux; external references with prefix #1 unavailable}%
  }%
}
\let\externaldocumentorig\externaldocument
\renewcommand{\externaldocument}[2][]{%
  \IfFileExists{tmp/xref/#2.aux}{%
    \externaldocumentorig[#1]{tmp/xref/#2}[#2.pdf]%
  }{%
  \IfFileExists{#2.aux}{%
    \externaldocumentorig[#1]{#2}%
  }{%
    \PackageWarning{stack}{Missing #2.aux; external references with prefix #1 unavailable}%
  }%
  }%
}

% Package for hypertext links:
\usepackage{hyperref}

% For any local file, say "hello.tex" you want to link to please
% use \externaldocument[hello-]{hello}
\externaldocument[introduction-]{introduction}
\externaldocument[conventions-]{conventions}
\externaldocument[sets-]{sets}
\externaldocument[categories-]{categories}
\externaldocument[topology-]{topology}
\externaldocument[sheaves-]{sheaves}
\externaldocument[sites-]{sites}
\externaldocument[stacks-]{stacks}
\externaldocument[fields-]{fields}
\externaldocument[algebra-]{algebra}
\externaldocument[brauer-]{brauer}
\externaldocument[homology-]{homology}
\externaldocument[derived-]{derived}
\externaldocument[simplicial-]{simplicial}
\externaldocument[more-algebra-]{more-algebra}
\externaldocument[smoothing-]{smoothing}
\externaldocument[modules-]{modules}
\externaldocument[sites-modules-]{sites-modules}
\externaldocument[injectives-]{injectives}
\externaldocument[cohomology-]{cohomology}
\externaldocument[sites-cohomology-]{sites-cohomology}
\externaldocument[dga-]{dga}
\externaldocument[dpa-]{dpa}
\externaldocument[sdga-]{sdga}
\externaldocument[hypercovering-]{hypercovering}
\externaldocument[schemes-]{schemes}
\externaldocument[constructions-]{constructions}
\externaldocument[properties-]{properties}
\externaldocument[morphisms-]{morphisms}
\externaldocument[coherent-]{coherent}
\externaldocument[divisors-]{divisors}
\externaldocument[limits-]{limits}
\externaldocument[varieties-]{varieties}
\externaldocument[topologies-]{topologies}
\externaldocument[descent-]{descent}
\externaldocument[perfect-]{perfect}
\externaldocument[more-morphisms-]{more-morphisms}
\externaldocument[flat-]{flat}
\externaldocument[groupoids-]{groupoids}
\externaldocument[more-groupoids-]{more-groupoids}
\externaldocument[etale-]{etale}
\externaldocument[chow-]{chow}
\externaldocument[intersection-]{intersection}
\externaldocument[pic-]{pic}
\externaldocument[weil-]{weil}
\externaldocument[adequate-]{adequate}
\externaldocument[dualizing-]{dualizing}
\externaldocument[duality-]{duality}
\externaldocument[discriminant-]{discriminant}
\externaldocument[derham-]{derham}
\externaldocument[local-cohomology-]{local-cohomology}
\externaldocument[algebraization-]{algebraization}
\externaldocument[curves-]{curves}
\externaldocument[resolve-]{resolve}
\externaldocument[models-]{models}
\externaldocument[functors-]{functors}
\externaldocument[equiv-]{equiv}
\externaldocument[pione-]{pione}
\externaldocument[etale-cohomology-]{etale-cohomology}
\externaldocument[proetale-]{proetale}
\externaldocument[relative-cycles-]{relative-cycles}
\externaldocument[more-etale-]{more-etale}
\externaldocument[trace-]{trace}
\externaldocument[crystalline-]{crystalline}
\externaldocument[spaces-]{spaces}
\externaldocument[spaces-properties-]{spaces-properties}
\externaldocument[spaces-morphisms-]{spaces-morphisms}
\externaldocument[decent-spaces-]{decent-spaces}
\externaldocument[spaces-cohomology-]{spaces-cohomology}
\externaldocument[spaces-limits-]{spaces-limits}
\externaldocument[spaces-divisors-]{spaces-divisors}
\externaldocument[spaces-over-fields-]{spaces-over-fields}
\externaldocument[spaces-topologies-]{spaces-topologies}
\externaldocument[spaces-descent-]{spaces-descent}
\externaldocument[spaces-perfect-]{spaces-perfect}
\externaldocument[spaces-more-morphisms-]{spaces-more-morphisms}
\externaldocument[spaces-flat-]{spaces-flat}
\externaldocument[spaces-groupoids-]{spaces-groupoids}
\externaldocument[spaces-more-groupoids-]{spaces-more-groupoids}
\externaldocument[bootstrap-]{bootstrap}
\externaldocument[spaces-pushouts-]{spaces-pushouts}
\externaldocument[spaces-chow-]{spaces-chow}
\externaldocument[groupoids-quotients-]{groupoids-quotients}
\externaldocument[spaces-more-cohomology-]{spaces-more-cohomology}
\externaldocument[spaces-simplicial-]{spaces-simplicial}
\externaldocument[spaces-duality-]{spaces-duality}
\externaldocument[formal-spaces-]{formal-spaces}
\externaldocument[restricted-]{restricted}
\externaldocument[spaces-resolve-]{spaces-resolve}
\externaldocument[formal-defos-]{formal-defos}
\externaldocument[defos-]{defos}
\externaldocument[cotangent-]{cotangent}
\externaldocument[examples-defos-]{examples-defos}
\externaldocument[algebraic-]{algebraic}
\externaldocument[examples-stacks-]{examples-stacks}
\externaldocument[stacks-sheaves-]{stacks-sheaves}
\externaldocument[criteria-]{criteria}
\externaldocument[artin-]{artin}
\externaldocument[quot-]{quot}
\externaldocument[stacks-properties-]{stacks-properties}
\externaldocument[stacks-morphisms-]{stacks-morphisms}
\externaldocument[stacks-limits-]{stacks-limits}
\externaldocument[stacks-cohomology-]{stacks-cohomology}
\externaldocument[stacks-perfect-]{stacks-perfect}
\externaldocument[stacks-introduction-]{stacks-introduction}
\externaldocument[stacks-more-morphisms-]{stacks-more-morphisms}
\externaldocument[stacks-geometry-]{stacks-geometry}
\externaldocument[moduli-]{moduli}
\externaldocument[moduli-curves-]{moduli-curves}
\externaldocument[examples-]{examples}
\externaldocument[exercises-]{exercises}
\externaldocument[guide-]{guide}
\externaldocument[desirables-]{desirables}
\externaldocument[coding-]{coding}
\externaldocument[obsolete-]{obsolete}
\externaldocument[fdl-]{fdl}
\externaldocument[index-]{index}

% Heap Project documents
\externaldocument[k3-]{k3}
\externaldocument[torelli-]{torelli}
\externaldocument[linear-algebra-]{linear-algebra}
\externaldocument[groups-]{groups}
\externaldocument[complex-analysis-]{complex-analysis}
\externaldocument[functional-analysis-]{functional-analysis}
\externaldocument[categories-]{categories}
\externaldocument[commutative-algebra-]{commutative-algebra}
\externaldocument[ringed-spaces-]{ringed-spaces}
\externaldocument[homological-algebra-]{homological-algebra}
\externaldocument[lie-algebras-]{lie-algebras}
\externaldocument[representation-theory-]{representation-theory}
\externaldocument[smooth-manifolds-]{smooth-manifolds}
\externaldocument[differential-topology-]{differential-topology}
\externaldocument[algebraic-topology-]{algebraic-topology}
\externaldocument[differential-geometry-]{differential-geometry}
\externaldocument[algebraic-geometry-]{algebraic-geometry}
\externaldocument[symplectic-geometry-]{symplectic-geometry}
\externaldocument[complex-geometry-]{complex-geometry}
\externaldocument[hodge-theory-]{hodge-theory}
\externaldocument[%
variations-of-hodge-structures-]{%
variations-of-hodge-structures}
\externaldocument[%
homological-mirror-symmetry-]{%
homological-mirror-symmetry}
\maybeexternaldocument{deformations-}{deformations}
\externaldocument[vertex-algebras-]{vertex-algebras}
\externaldocument[springer-]{springer}
\externaldocument[infty-categories-]{infty-categories}
\externaldocument[seminars-differential-geometry-]{%
seminars-differential-geometry}
\externaldocument[seminars-complex-geometry-]{%
seminars-complex-geometry}
\externaldocument[seminars-algebraic-geometry-]{%
seminars-algebraic-geometry}
\externaldocument[seminars-analysis-pde-]{%
seminars-analysis-pde}
\externaldocument[surfaces-]{surfaces}
\externaldocument[birational-maps-commutative-algebra-]{birational-maps-commutative-algebra}
\externaldocument[cremona-transformations-]{cremona-transformations}
\externaldocument[derived-categories-of-sheaves-]{derived-categories-of-sheaves}
\externaldocument[geometric-invariant-theory-]{geometric-invariant-theory}
\externaldocument[geometric-stability-conditions-and-group-actions-]{geometric-stability-conditions-and-group-actions}
\externaldocument[moduli-spaces-of-sheaves-]{moduli-spaces-of-sheaves}
\externaldocument[bridgeland-stability-general-theory-]{bridgeland-stability-general-theory}
\externaldocument[hilbert-schemes-]{hilbert-schemes}
\externaldocument[noncommutative-abelian-surfaces-and-kummer-type-hyperkahler-varieties-]{noncommutative-abelian-surfaces-and-kummer-type-hyperkahler-varieties}
\externaldocument[mumford-tate-groups-in-hodge-theory-]{mumford-tate-groups-in-hodge-theory}
\externaldocument[hodge-birational-atoms-2026-]{hodge-birational-atoms-2026}
\externaldocument[xxii-egd-2026-]{%
xxii-egd-2026}
\externaldocument[pde-]{pde}
\externaldocument[probability-]{probability}

\externaldocument[linguistics-]{linguistics}

\externaldocument[%
16th-alga-meeting-2026-]{%
16th-alga-meeting-2026}

\externaldocument[%
iv-jornadas-lefschetz-2026-]{%
iv-jornadas-lefschetz-2026}

\externaldocument[reading-the-masters-]{%
reading-the-masters}

\externaldocument[real-analysis-]{%
real-analysis}

\externaldocument[optimal-transport-]{%
optimal-transport}

\externaldocument[history-]{history}

\externaldocument[seiberg-witten-]{seiberg-witten}

\externaldocument[history-philosophy-biology-]{history-philosophy-biology}

\externaldocument[classical-mechanics-]{classical-mechanics}
\externaldocument[electromagnetism-]{electromagnetism}
\externaldocument[waves-and-optics-]{waves-and-optics}
\externaldocument[wif-]{waves-in-fluids-iv}
\externaldocument[thermal-physics-]{thermal-physics}
\externaldocument[quantum-mechanics-]{quantum-mechanics}
\externaldocument[special-relativity-]{special-relativity}
\externaldocument[physics-]{physics}

\externaldocument[chemistry-]{chemistry}

\externaldocument[molecular-cell-biology-]{molecular-cell-biology}

\externaldocument[microbiology-]{microbiology}

\externaldocument[protistology-]{protistology}

\externaldocument[mycology-]{mycology}

\externaldocument[botany-]{botany}

\externaldocument[zoology-]{zoology}

\externaldocument[animal-histology-]{animal-histology}

\externaldocument[animal-physiology-]{animal-physiology}

\externaldocument[systematics-biogeography-]{systematics-biogeography}

\externaldocument[ecology-natural-resources-]{ecology-natural-resources}

\externaldocument[earth-sciences-]{earth-sciences}

\externaldocument[biostatistics-]{biostatistics}

\externaldocument[mathematical-systems-biology-]{mathematical-systems-biology}

\externaldocument[cimpa-el-salvador-2024-]{cimpa-el-salvador-2024}

\externaldocument[anthropology-]{anthropology}

\externaldocument[yang-mills-theory-]{yang-mills-theory}

\externaldocument[seminars-algebra-]{seminars-algebra}

\externaldocument[seminars-physics-]{seminars-physics}

\externaldocument[seminars-computing-]{seminars-computing}

\externaldocument[oktobermat-2026-]{oktobermat-2026}

\externaldocument[saet-2026-]{saet-2026}

\externaldocument[seminars-symplectic-geometry-]{seminars-symplectic-geometry}

\externaldocument[seminars-dynamics-]{seminars-dynamics}

\externaldocument[oktobermat-2024-]{oktobermat-2024}

\externaldocument[quantum-groups-]{quantum-groups}

\externaldocument[legal-2027-]{legal-2027}

\externaldocument[derived-categories-]{derived-categories}

% Theorem environments.
%
\theoremstyle{plain}
\newtheorem{theorem}[subsection]{Theorem}
\newtheorem{proposition}[subsection]{Proposition}
\newtheorem{lemma}[subsection]{Lemma}
\newtheorem{corollary}[subsection]{Corollary}

\theoremstyle{definition}
\newtheorem{definition}[subsection]{Definition}
\newtheorem{example}[subsection]{Example}
\newtheorem{exercise}[subsection]{Exercise}
\newtheorem{situation}[subsection]{Situation}

\theoremstyle{remark}
\newtheorem{remark}[subsection]{Remark}
\newtheorem{remarks}[subsection]{Remarks}

\numberwithin{equation}{subsection}

% Macros
%
\def\lim{\mathop{\mathrm{lim}}\nolimits}
\def\colim{\mathop{\mathrm{colim}}\nolimits}
\def\Spec{\mathop{\mathrm{Spec}}}
\def\Hom{\mathop{\mathrm{Hom}}\nolimits}
\def\Ext{\mathop{\mathrm{Ext}}\nolimits}
\def\SheafHom{\mathop{\mathcal{H}\!\mathit{om}}\nolimits}
\def\SheafExt{\mathop{\mathcal{E}\!\mathit{xt}}\nolimits}
\def\Sch{\mathit{Sch}}
\def\Mor{\mathop{\mathrm{Mor}}\nolimits}
\def\Ob{\mathop{\mathrm{Ob}}\nolimits}
\def\Sh{\mathop{\mathit{Sh}}\nolimits}
\def\NL{\mathop{N\!L}\nolimits}
\def\CH{\mathop{\mathrm{CH}}\nolimits}
\def\proetale{{pro\text{-}\acute{e}tale}}
\def\etale{{\acute{e}tale}}
\def\QCoh{\mathit{QCoh}}
\def\Ker{\mathop{\mathrm{Ker}}}
\def\Im{\mathop{\mathrm{Im}}}
\def\Coker{\mathop{\mathrm{Coker}}}
\def\Coim{\mathop{\mathrm{Coim}}}

% Boxtimes
%
\DeclareMathSymbol{\boxtimes}{\mathbin}{AMSa}{"02}

%
% Macros for moduli stacks/spaces
%
\def\QCohstack{\mathcal{QC}\!\mathit{oh}}
\def\Cohstack{\mathcal{C}\!\mathit{oh}}
\def\Spacesstack{\mathcal{S}\!\mathit{paces}}
\def\Quotfunctor{\mathrm{Quot}}
\def\Hilbfunctor{\mathrm{Hilb}}
\def\Curvesstack{\mathcal{C}\!\mathit{urves}}
\def\Polarizedstack{\mathcal{P}\!\mathit{olarized}}
\def\Complexesstack{\mathcal{C}\!\mathit{omplexes}}
% \Pic is the operator that assigns to X its picard group, usage \Pic(X)
% \Picardstack_{X/B} denotes the Picard stack of X over B
% \Picardfunctor_{X/B} denotes the Picard functor of X over B
\def\Pic{\mathop{\mathrm{Pic}}\nolimits}
\def\Picardstack{\mathcal{P}\!\mathit{ic}}
\def\Picardfunctor{\mathrm{Pic}}
\def\Deformationcategory{\mathcal{D}\!\mathit{ef}}


\begin{document}

\title{25th SAET Conference}
\maketitle

\phantomsection
\label{section-phantom}

\tableofcontents

25th Annual SAET Conference, held at
IMPA in Rio de Janeiro, Brazil, from July 13
to July 18, 2026.


\section{Scheinkman}
\label{section-scheinkman}

\noindent
The upshot of Scheinkman's lecture
at SAET is:
Can economic incentives be designed so that
preserving nature wins against destroying
it?

And apparently we expect his answer to be
affirmative: Properly priced carbon capture
can make forest recovery more profitable
than low-productivity cattle ranching.



\section{Growth and Redistribution with
Risk Lovers and Risk Averters}
\label{%
section-risk-lovers%
}

\medskip\noindent
Timothy Kehoe,
University of Minnesota.

\medskip\noindent
This talk was given at the 25th Annual
SAET Conference, held at IMPA in Rio de
Janeiro, Brazil, from July 13 to
July 18, 2026.

\medskip\noindent
{\bf Short communication.}

\medskip\noindent
{\bf Model.}
The economy consists of two types of
agents: risk lovers and risk
averters. Risk lovers invest in a
risky, high-return technology,
whereas risk averters prefer a safer
asset with lower returns.

\medskip\noindent
{\bf Main question.}
The model studies the trade-off
between economic growth and
redistribution. The risky investment
generates growth, but also increases
inequality.

\medskip\noindent
{\bf Welfare function.}
The objective is to maximize a CES
social welfare function. Different
choices of the CES parameter assign
different weights to inequality.

\medskip\noindent
{\bf Optimal policy.}
The model computes the optimal tax
rate. This is analogous to finding an
optimal policy parameter (such as the
optimal emission price in the forest
paper): the parameter maximizing the
chosen welfare function.

\medskip\noindent
{\bf Interpretation.}
Since the risk lovers become richer,
the policy effectively taxes the rich
and redistributes resources to the
risk averters. The interesting
question is how much redistribution
can be introduced before growth is
reduced too much.

\medskip\noindent
{\bf Remarks.}
The classification "risk lover" is a
technical one: it refers to a
preference for risky gambles, not to
a political or moral category.
Likewise, "optimal" always depends on
the chosen welfare function.


\section{Real Effects of Interest-Rate
Policy and Central-Bank Asset Purchases
with Endogenous Collateral}
\label{%
section-interest-rate-policy%
}

\medskip\noindent
Susan Schommer,
University of Bonn.

\medskip\noindent
This talk was given at the 25th Annual
SAET Conference, held at IMPA in Rio de
Janeiro, Brazil, from July 13 to
July 18, 2026.

\medskip\noindent
{\bf Upshot (human-written).}


\medskip\noindent
{\bf Motivation.}
In the basic flexible-price economy,
$p_1$ and $\sigma^h$ are
predetermined. This does not capture
the effects of predictable monetary
policy.

\medskip\noindent
{\bf Household problem.}
Given $(p,q,y)$, household $h$
chooses a vector in order to maximize
its utility.

\medskip\noindent
{\bf Equilibrium allocation.}
The equilibrium allocation is the
equilibrium vector obtained by
collecting the optimal choices of all
households (and firms). Its
coordinates correspond to quantities
such as consumption, investment,
asset holdings, etc.

\medskip\noindent
{\bf Theorem (equilibrium mapping).}
Every equilibrium of the basic
price-friction economy
$\mathcal E$ maps to an equilibrium
of the auxiliary flexible-price
economy $\widehat{\mathcal E}$.

\medskip\noindent
{\bf Proposition (real effects of
interest policy).}
Consider the basic price-friction
economy $\mathcal E$ with logarithmic
utility and fixed $x^{CB}_3$.
Changes in the policy interest rate
$i\ge0$ have real effects.
Specifically, variations in $i$
affect relative prices, real rates of
return, and the equilibrium
allocation.

\medskip\noindent
{\bf Example.}
One figure plots the policy interest
rate against several coordinates of
the equilibrium allocation. The
curves are not constant, showing that
different coordinates of the
equilibrium vector vary with $i$.

\medskip\noindent
{\bf Interpretation.}
The basic flexible-price model does
not detect real effects of monetary
policy. Introducing a single price
friction restores them: the model now
detects that changing the policy
interest rate changes the equilibrium
allocation.

\medskip\noindent
{\bf Extended model.}
The extended model incorporates
predictable monetary policy. The
conclusion is that only the
unforecastable component of monetary
policy has real effects, since
predictable changes are already taken
into account by the agents.

\medskip\noindent
{\bf Conventional versus
unconventional monetary policy.}
Conventional monetary policy changes
the short-term policy interest rate.
Unconventional monetary policy acts
through changes in the central-bank
balance sheet, for instance by
purchasing assets.

\medskip\noindent
{\bf Conclusions.}
The optimal interest-rate policy is
state-contingent: the optimal policy
depends on the state of the economy,
and interest rates should be raised
in high-production states.

For unconventional monetary policy,
the effects depend on the financial
constraint. Relaxing leverage
constraints improves welfare, whereas
relaxing short-sale constraints does
not.

\section{Fiscal Sustainability and
Macroeconomic Stabilization}
\label{section-fiscal-sustainability}

\medskip\noindent
Michael Woodford,
Columbia University.

\medskip\noindent
This talk was given at the 25th Annual
SAET Conference, held at IMPA in Rio de
Janeiro, Brazil, from July 13 to
July 18, 2026.

\medskip\noindent
{\bf Introduction.} Policy analysis involves
working with ``infinite-horizon general
equilibrium models'' --- somehow ideal
states of equilibrium predicted by the model
which would occur only at infinite time.

[Missing content.]

Somehow this might cause that certain
{\it transfer policies} have no
consequences.

Asumption: sustainability can be treated as
axiomatic or not. The choice might determine
what monetary policy you will choose
to stabilize marcoeconomics.

Consider monetary policy rules of
Taylor-type form:
$$
i_t=\phi_\pi \pi_t + \phi_y y_t+
e^m_t,
$$
with $\phi_\pi,\phi_y \geq  0$,
$\varepsilon_t^m$ exogenous (may respond to
other ``fundamental'' disturbances),
and all variables written in terms of
deviations from a zero-inflation steady
state.

Then in a standard RANK model,
log-linearized equations can be writen in
vector form [involves a matrix].

\medskip\noindent
{\bf Recap of REE analysis.}
(Rational expectations ?.)

Idea: no feedback to nominal interest rates
from (?) inflation.

\medskip\noindent
{\bf Adding debt dynamics.}
Suppose that the only debt instrument is a
one-period nominal bond, and that all
taxes/transfers are lump-sum.

Then previous structural equations are
unchanged; but we must add an equation
describing the dynamics of public debt,
taking into account transfer policy:
[equation…]

\medskip\noindent
{\bf }
Hence REE implies
\begin{itemize}
\item AM/PF or AF/PM $\implies $ unique
non-explosive REE.
\item AM/AF $\implies$ only explosive
solutions.
\item PM/PF $\implies$ indereminacy, sunspot
equation.
\end{itemize}

\noindent
In this analysis, it does matter
whether there is a commitment to sustainable
fiscal policy.

[We're looking at how to find a proper
macroeconomics stabilization policy.]

\medskip\noindent
{\bf Questioning this analysis.}
\begin{itemize}
\item Conceptual problem: how can one settle
whether fiscal policy guarantees
sustainability or no, when this is entirely
a question about what fiscal policy promises
to do infinitely far in the future?

\item (Missing)
\end{itemize}

\medskip\noindent
{\bf Alternative approach.}

\begin{itemize}
\item Can we model the way in which people
make forward-looking decisions, without
assuming (unrealistically) that they have
definite expectations about the far
future…

\item  The idea: 
\begin{itemize}
\item Economic agents formulate plans some
finite distance into the future, on the
basis of a model of possible transitions
over that planning horizon (don't look at
all possible final states, but a reduced
set).

\item the stages that can be reached at the
horizon at which forward planning is
truncated are evaluated using a {\it learned
value function}.
\item agents don't continue to execute plans
made in the past, but re-optimize from
whatever situation they have reached.
\item value functions are update as
additional experience is accumulated,
through reinforcement learning.
(AAAAAA!!!!)
\end{itemize}
\end{itemize}

\medskip\noindent
{\bf Learning the value functions.}
\begin{itemize}
\item As part of its planning exercise in
period $t$, each houshold $h$ computes an
estimate of the value of its objective
(expected discounted utility from $t$
onward), for any (counterfactual) level of
wealth $b_t^h = b$ that is might have
brought into the period --- call this
function $v^{est}_t(b)$.

\item Each firm $f$ similarly…
\item Value function can then be updated for
the next period's exercise using error
correction rule (difference equation).
\end{itemize}


\medskip\noindent
{\bf Rational expectations?}
\begin{itemize}
\item What relationship is ther ebetween the
dynamics implied by a computational model of
this kind and REE predictions?

\item If the value functions used by
households and firms in period $t$ are the
ones implied by a particular REE they will
all take the actions associate dwith the
REE.

\item Quesion: whether learning dynamics
converge? …?
\end{itemize}

\medskip\noindent
{\bf Log-Linear leaning dynamics.}
\begin{itemize}
\item The solution to this linear system
then determins
$$
b_{t+1}=b_{t+1|_t},\qquad 
v_t^{\operatorname{est}}=x_{t|t}-
\zeta b_t
$$
and …
if so the value functions converge…

\medskip\noindent
{\bf When do learning dynamics converge?}
\begin{itemize}
\item Conclusion when
$\bar{b}=0$:convergence if and only if AM/PF
policy configuration --- either passive MP
or active FP produces an unstable
eigenvalue.

\item Continues to be true for (modestly)
values of…
\end{itemize}

\medskip\noindent
{\bf Numerical example: $\bar{b}=1.0$ GDP, 5
year horizon.} (Figure: AM/PF remains
sufficient, and almost necessary for
stability.)

\medskip\noindent
{\bf Numerical example: finer detail.}
(Similar figure.)

\medskip\noindent
{\bf Conclusions.}
\begin{itemize}
\item One shouldn't expect transfer policy
to be irrelevant to aggregate demand
determination (and hence for price
stability)

\item According to our model with
finite-horizon planning…

\item Broader lesson: when analysing
possible stabilization policies we
shouldn't be content merely to consider
the REE consistent with a given policy
commitment we should also ask whether ar
least something close to that outcome can be
achieved with a more realistic model of
peoples thought processes --- not always the
case!
\end{itemize}


\end{itemize}

\medskip\noindent
{\bf Reference.}
Ayres et al.\ (2019) study the
monetary and fiscal history of Brazil
from 1960 to 2016, especially the
hyperinflation episodes.




\section{Inflation Targeting under
Fiscal Fragility}
\label{section-inflation-targeting}

\medskip\noindent
Vitor Chagas da Costa,
IMPA.

\medskip\noindent
This talk was given at the 25th Annual
SAET Conference, held at IMPA in Rio de
Janeiro, Brazil, from July 13 to
July 18, 2026.


\medskip\noindent
{\bf Question.} How to set the inflation
targe in an economy with high government
debt burden?

\begin{itemize}
\item Some authority has an inflation target
and attempts to implement policies to
achieve it. But in practice, the target is
not attained.

\item High debt burden raises the temptation
of the policymaker to use inflation as a
partial default mechanism. However, high
inflation reduces welfare and affects
monetary policy credibility.

\item (missing.)

\item Upshot of the paper:
``The target is important in balancing
this tradeoff''
\end{itemize}

\medskip\noindent
{\bf Brazilian experience with inflation
targeting regime.} (graph: effective
inflation and inflation target 
(the latter with a neighbourhood: lower
and upper limit). The actual data is much
more random and often outside the expected
neighbourhood!)

\medskip\noindent
{\bf Model.} We model an economy with two
types of agents:
\begin{itemize}
\item An altruistic policymaker acts as a
fiscal authority and central bank, under an
inflation target regime.

\item Private agents form rational
expectations and hold debt.
\item The PM chooses the combination of
monetary and fiscal policy that maximizes
welfare, given an initial debt level
\end{itemize}

\noindent
Deviation from target solution (…)

\medskip\noindent
{\bf Results.} What the policy would do is
dependent on that limit:

\begin{itemize}
\item we obtain endogenous limits that determine
target coordination ability.

\begin{itemize}
\item fiscal dominance zone: high debt,
target irrelevant
\item no crisis zone: low debt, target
delivered.
\item in between: fiscal fragility
\end{itemize}

\item In the fiscal fragility zone (like in
Brazil's economy!):
\begin{itemize}
\item expected inflation is higher than
target.
\item delivering the target is more costly.
We need high interest rate and physical
effort to attain the target.
\end{itemize}

\item Raising inflation target:
\begin{itemize}
\item reduces the real interest rate.
\item increases lower bound of FFZ.
\item reduces the time to exit the FFZ.
\end{itemize}

\item The optimal target is the lowest
possible (inflation level) such that 
it is credible.
\end{itemize}

\medskip\noindent
{\bf References}Basic to inflation: Ser and
Wallace…

\medskip\noindent
{\bf Introducing the model: policymaker}
\begin{itemize}
\item Chooses inflation $\pi_t \geq 0$
and debt $D_{t+1}$
$$
\operatorname{max}_{\pi_t,D_{t+1}}
\mathbb{E}_0 \sum_{t=0}^\infty 
\beta^t[u(c_t,g_t)-\psi(\pi_t)]
$$
subject to $g_t+(1+r_t) D_t \leq 
D_{t+1}+\tau(e-\varepsilon_t)$.

We assume:
\begin{itemize}
\item $c_t,g_t>0$
\item $u(c,g)=\rho_c+(1-\rho)v(g)$, utility
in private consuption is risk-neutral,
and $v$ is strictly concave.
\item $\psi$ is the disutility of inflation,
such that $\psi'>0$ and $\psi'' \geq 0$.
\item Endowment $e$ taxed at fixed rate $\tau$
\item $\varepsilon_t$ is the endowment
penalty for deviating from target.
\end{itemize}
\end{itemize}

\medskip\noindent
{\bf Private agents.} (Missing, formulas and
parameters.)


\medskip\noindent
{\bf Policy maker's choices.}
The first-order condition for debt from the
private agents problem implies the following
equilibrium of real interest rate:
(formula).

\medskip\noindent
{\bf Timing.} Interest rate is defined in
equilibrium.

\begin{enumerate}
\item Sunspot variable $\zeta_t$ is
realized.
\item Policymakes chooses to follow the
target or not, given realization of
$\zeta_t$.
\item …
\item …
\end{enumerate}

\begin{itemize}
\item If expected inflation is high,
by the equation, the policymaker may decide
to inflate the debt away, so, it's costly.
\end{itemize}

\medskip\noindent
{\bf Fisical Fragility.}
As a result, we find, as we already
anticipated, there are three possible zones
(which depend on belief of the agents):

\begin{itemize}
\item (No crisis zone.) 
Agents believe the goverment will
deviatie: it's optimal for the goverment to
follow the target. This is the only solution
of the model.

\item (Fiscal dominance zone.) Even if agents believe the goverment,
the government might deviate.

\item (Fiscal fragility zone.) If agents believe believe the
goverment wants to deviate (not keep
inflation on target),
\end{itemize}

\noindent
In FEZ, estimated inflation is given by
[formula].

\medskip\noindent
{\bf Main result of paper.}
Higher target improves coordination by
reducing inflation volatility.
Therefore, the rollover cost is reduced,
lowering the temptation of deviating.

\medskip\noindent
{\bf Next result.}
The boundary between no crisis and fiscal
fragility…

\medskip\noindent
{\bf Numerical example.}
In 2002 Brazilian confidence crisis.
Inflation expectations grew by
fear that the president would change
policies.

\noindent
{\bf Optimal fiscal policy.}
Figure: in fiscal fragility, the policymaker
reduces (?) to get back to no crisis zone.

We can also plot the same functions varying
the target.

\medskip\noindent
{\bf Concluding remarks.}
\begin{itemize}
\item A higher inflation target may ssist
policymaker in economies with high debt
burden that are subject to sudden loss
of confidence on the commitment with the
target.

\item Decreasing debt levels through
austerity serves as a fiscal support to
better target coordination when the economy
is in the lower part of the FFZ.

\item Model can be relevant to the current
experience with high inflation and high debt
in emergent and advanced economies (e.g.
USA, European economies).
\end{itemize}

\medskip\noindent
{\bf Audience questions:} 

There is only
one target in the model, but the whole thing
is changing the target. How to decide?


Did you consider
price-inflation targeting, average-inflation
targeting?

\section{Recursive Contracts in
Non-Convex Environments}
\label{section-recursive-contracts}

\medskip\noindent
Felix Kübler,
University of Zurich.

\medskip\noindent
This talk was given at the 25th Annual
SAET Conference, held at IMPA in Rio de
Janeiro, Brazil, from July 13 to
July 18, 2026.

\medskip\noindent
{\bf Non-convexities in Aloisio's work.}
\begin{itemize}
\item Asymmetric information, 2001.
\item Non-convex budget sets, 2002.
\item Risk-loving preferences, 2018.
\end{itemize}


\medskip\noindent
{\bf Lotteries.}
\begin{itemize}
\item Many economic models involve
non-convex optimization problems.
\item Non-convexities.
\begin{enumerate}
\item Difficult to solve optimization
problems with nonconvex constraints or
objective function that is not concave.
\item Lotteries can increase the value of
the objective function.
\end{enumerate}
\item Since allowing for randomization
across allocations might increase the value
of the objective function, there has been
substantial interest in exploring lottery
solutions to models with nonconvexities
(Weiss 1976, Myerson 1982, Prescott and
Townsend 1984, Arnott-Stiglitz 1988.)

\item Questions:
\begin{enumerate}
\item Does it make sense?
\item How can we compute optimal loteries?
\end{enumerate}
\end{itemize}

\medskip\noindent
{\bf Simple example} Maximize $f$ with
constraints $f \leq  A$, $f \leq B$.

\medskip\noindent
{\bf Ramsey problem (economic example).} 
Consider the following non-convex problem:
(maximization, constraints, $T$ representing
taxes…)
The housholds know the probability
distribution…

Do lotteries ``make sense''?
Yes - if lottery is ex post (see Citanna and
Tvede 2024 and our paper for examples).


\medskip\noindent
{\bf Dynamic problem with forward looking
constraint.}

\medskip\noindent
{\bf Problems.} MM 2019 assumes concave
objective + convex constraint set. Many
important applications violate this: e.g.
Ramsey problem. Prior results (Cole-Kubler
2012)…

\medskip\noindent
{\bf Several slides missing.}

\medskip\noindent
{\bf First period.} Involving
sub-gradient descent.

\medskip\noindent
{\bf Algorithm for computing optimal
lotteries (Algorithm 1).}

\noindent
{\bf Given:} value function $D(\cdot)$,
state $(x,s)$, initial promise
$\phi$, initial multipliers
$\lambda^1, \mu^1$.

\noindent
{\bf Subgradient Descent on Multipliers.} 
For $k=1,…,N$:
\begin{enumerate}
\item Solve $a^k \in \operatorname{arg}
\operatorname{max}_a[$ long formula].
\item Pick $\phi^k(s') \in \partial
\partial D(\lambda^k+\mu^k,x',s')$.
\item Update.
\end{enumerate}

\noindent
{\bf Output.} Stage lottery.

\medskip\noindent
{\bf Slightly more interesting example.}

\medskip\noindent
{\bf Conclusion.}
\begin{itemize}
\item Equivalence between the dual and the
lottery relaxation.
\item Different recusive formulas
reflect different versions of lottery
relaxation.
\item Subgradient descend method
to recover the policy.
\item Application to Ramsey problem.
\end{itemize}

\section{From Degradation to Recovery:
Quantifying the Impact of Emission
Prices on Tropical Forests}
\label{section-tropical-forests-scheinkman}

\medskip\noindent
José Scheinkman,
Columbia University.

\medskip\noindent
This talk was given at the 25th Annual
SAET Conference, held at IMPA in Rio de
Janeiro, Brazil, from July 13 to
July 18, 2026.

\medskip\noindent
{\bf Motivation.}
Want to quantify impact of emission prices
on tropical rainforests.

\begin{itemize}
\item The 
Brazilian Amazon provided the best data set,
60% of the whole Amazon. Also, has been
substatianlly deforested.

\item Many scientists predicted parts of the
Amazon were closed to being ?.

\item There's a lot of work already on the
Amazon. Bolsonaro = lots of deforestation.
This year = least deforestation.

\item Deforestation is a dynamical process.
\end{itemize}

\medskip\noindent
{\bf Today I.}

If stays the same: emit $\approx 14.5$ GT
of $CO_2$ in 30 years.

The second part: effect of this in
biodiversity, is work in progress.

\medskip\noindent
{\bf State and control variables.}
\begin{itemize}
\item $Z$…
\item $P^a$ will be a constant for today.
\end{itemize}

\begin{remark}
\label{remark-carbon-cattle}
There's much better feed for carbon
than for cattle.
\end{remark}

\medskip\noindent
{\bf Robust optimization II.}
There are ``penaliazations'' (this reminds
me of RL): penalize departure.

\medskip\noindent
{\bf Planner's shadow price I
(or the previous slide).}
Use Hamiltonian Markov chain.

We check if we can predict the
information from the past!

The only tool you have is how much you value
deforestation, so you adjust the model.

\medskip\noindent
{\bf Before "biodiversity I"}
Somehow handle income/interest, money
basically, so that deforestation is no
longer convenient (maybe the whole point,
right?).

\medskip\noindent
{\bf biodiversity I.}
Arrhenius 1921 method works great
for one or two hectares.

\medskip\noindent
{\bf biodiversity II.}
Model: how many species that have n many
individuals in x hectares? Fischer's $\alpha$
determines the average number of species.
People change the $\alpha$ as a function
of the area.

\medskip\noindent
{\bf Clustering.}
Species are not evenly distributed,
so we need clustering. Clustering makes
i
\medskip\noindent
{\bf Forest cover changes…}

\medskip\noindent
{\bf There's not a conclusion slide but…}
Lesson: there's a lot of potential
for reforesting the Amazon.
People are not being paid for recuperating
the Amazon. All we do is penalize CO2
emissions. You get your biodiversity pack…
They have proposed, negotiated, to get this
done. The basis is this paper.

In Brazil we have data at municipal level
every 5 years.

Nobody knows how to value biodiversity.
``I'm gonna restore biodiversity:
give me $x$ dollars for every… what
restored?''
Biodiversity is a very complex system
--- a very complicated dynamical system.
This happens in any rainforest.
We don't understand the dynamics of
tropical rainforests.

Who does the reforestation: the government
could compensate organizations, but
most likely not doing it itself.
The average hectare has 5000 tons of CO2, 
4000 dollars. That is: here's 4000 dollars,
reforest the land.

30 years is enough to restore the Amazon ---
in 30 years you are in the 90s.
That number $\alpha$ is/takes the convergence
rate,

Aloiso: it's a very big economic
problem: who's gonna pay for this?


S: someone says:
 I'm gonna pay for this --- then someone
else will deforest… what's the use?
There are negotiations, like the ETS,
Brazil-China to make this happen.
China communist party.

\bibliography{my}
\bibliographystyle{amsalpha}

\end{document}
